\documentclass[11pt, a4paper]{article}

\usepackage{amsmath, amssymb, amsthm}
\usepackage{geometry}
\usepackage{booktabs}
\usepackage{array}
\usepackage{xcolor}
\usepackage{hyperref}
\usepackage{microtype}
\usepackage{parskip}
\usepackage{enumitem}
\usepackage{titlesec}
\usepackage{fancyhdr}

\geometry{margin=2.5cm}

% Colours
\definecolor{clockblue}{RGB}{30, 80, 160}
\definecolor{syncgreen}{RGB}{20, 120, 60}
\definecolor{warmorange}{RGB}{190, 90, 20}

% Section formatting
\titleformat{\section}{\large\bfseries\color{clockblue}}{}{0em}{}[\titlerule]
\titleformat{\subsection}{\normalsize\bfseries\color{syncgreen}}{}{0em}{}

% Header / footer
\pagestyle{fancy}
\fancyhf{}
\rhead{\small\textit{Bell Without Faster Than Light}}
\lhead{\small\textit{Time-Phase Clock Equations}}
\cfoot{\thepage}

% Boxed equation environment
\newcommand{\beq}[1]{%
  \begin{center}
  \fbox{\begin{minipage}{0.88\textwidth}\vspace{4pt}%
  \begin{equation}#1\end{equation}%
  \vspace{0pt}\end{minipage}}%
  \end{center}}

\title{
  {\Large\bfseries Bell Without Faster Than Light}\\[6pt]
  {\large Equations of the Time-Phase Clock Interpretation}\\[4pt]
  {\normalsize\itshape A Radiologist's Toy Theory}
}
\author{rayolddog \quad \texttt{github.com/rayolddog/BellWithoutFasterThanLight}}
\date{}

\begin{document}

\maketitle
\thispagestyle{fancy}

\begin{abstract}
The phase of the Schr\"odinger wave function is interpreted as an internal clock carried
by each particle, oscillating at $\omega = E/\hbar$.  Entanglement is synchronized clocks
at creation.  Measurement is Kuramoto synchronization with a local detector.
The wave never collapses.  No faster-than-light signal is required.
This document collects the key equations developed across the project,
from the Bell correlation through Dirac spinors, the Higgs mechanism,
gravity, Penrose objective reduction, twistor theory, thermodynamics,
Brownian motion, and the Heisenberg uncertainty principle.
\end{abstract}

\tableofcontents
\vspace{1em}
\hrule
\vspace{1em}

% ─── 1 ────────────────────────────────────────────────────────────────────────
\section{The Time-Phase Wave Function}

The Schr\"odinger equation -- which describes orbitals exactly and never needs to collapse:
\begin{equation}
  i\hbar \frac{\partial \psi}{\partial t} = \hat{H}\psi
\end{equation}

The wave function augmented with an explicit internal clock phase $\varphi$:
\begin{equation}
  \psi_A(\mathbf{x}, t)
  = \sum_s c_s\,|s\rangle \cdot e^{-i\omega t + i\varphi_A(t)}
\end{equation}
where $\omega = E/\hbar$ is the particle's clock frequency and $\varphi_A(t)$ is its
internal clock phase, distinct from the dynamical phase $-i\omega t$.

The probability of detecting a property along axis $\hat{n}$ depends on the
\textbf{relative phase} between particle clock $\varphi_A$ and detector clock $\Phi_D$:
\begin{equation}
  P\!\left(+1\,\middle|\,\hat{n},\,\varphi_A,\,\Phi_D\right)
  = \cos^2\!\left(\frac{\varphi_A - \Phi_D - \theta_{\hat{n}}}{2}\right)
\end{equation}

\textbf{Entanglement} is synchronized clocks at the moment of creation:
\begin{equation}
  \varphi_A(0) = \varphi_B(0) = \varphi_0
  \qquad\text{(shared initial phase, one spacetime point)}
\end{equation}

% ─── 2 ────────────────────────────────────────────────────────────────────────
\section{Kuramoto Phase Synchronization}

The Kuramoto (1975) equation governs clock synchronization and thereby measurement:
\beq{
  \frac{d\varphi}{dt} = \omega + K\sin(\Phi_{\text{bulk}} - \varphi)
  \label{eq:kuramoto}
}

\begin{itemize}[nosep]
  \item Free flight: $K = 0 \;\Rightarrow\; \varphi(t) = \varphi_0 + \omega t$
  \item At detector: $K > 0 \;\Rightarrow\; \varphi \to \Phi_{\text{bulk}}$ \;(synchronization)
\end{itemize}

Synchronization fixed point: $\varphi^\star = \Phi_{\text{bulk}}$.
Synchronization timescale: $\tau_{\text{sync}} = 1/K$.

Both detectors are pre-synchronized by thermal equilibration with the macroscopic environment:
\begin{equation}
  \Phi_{D_1} = \Phi_{D_2} = \Phi_{\text{bulk}}
  \qquad\text{(established locally, before the experiment)}
\end{equation}

% ─── 3 ────────────────────────────────────────────────────────────────────────
\section{Bell's Theorem -- Where the Factorization Fails}

Bell's hidden variable theorem:
\begin{equation}
  E(a,b) = \int A(a,\lambda)\,B(b,\lambda)\,\rho(\lambda)\,d\lambda
  \label{eq:bell}
\end{equation}

The hidden state in this model is $\lambda = (\varphi_0,\Phi_{\text{bulk}})$.
Because $\Phi_{\text{bulk}}$ is the \emph{same} variable at both detectors, the joint distribution
does not factorize:
\begin{equation}
  \rho(\varphi_0, \Phi_{\text{bulk}}) \;\neq\; \rho(\varphi_0)\cdot\rho(\Phi_{\text{bulk}})
\end{equation}
Bell's integral~\eqref{eq:bell} implicitly assumes independent components.
This is not a conspiracy: $\Phi_{\text{bulk}}$ is a single shared variable established
subliminally and locally, long before the measurement.

\begin{center}
\begin{tabular}{lp{4cm}p{5cm}}
\toprule
Assumption & Standard form & Status in this model \\
\midrule
Locality & $A$ depends only on $(a,\lambda)$ & Satisfied --- each detector is local \\
Measurement independence & $\rho(\lambda\mid a,b)=\rho(\lambda)$ & \textbf{Fails} --- $\Phi_{\text{bulk}}\in\lambda$ is shared \\
\bottomrule
\end{tabular}
\end{center}

\subsection{Correlation functions}

\begin{center}
\begin{tabular}{lcc}
\toprule
Model & $E(a,b)$ & $\text{CHSH}_{\max}$ \\
\midrule
Classical LHV (deterministic)    & $-\bigl(1 - 2|\Delta|/\pi\bigr)$ & $2.000$ \\
Phase-clock (Malus law)          & $-\tfrac{1}{2}\cos\Delta$        & $\sqrt{2}\approx 1.414$ \\
Standard quantum mechanics       & $-\cos\Delta$                    & $2\sqrt{2}\approx 2.828$ \\
\bottomrule
\end{tabular}
\end{center}

% ─── 4 ────────────────────────────────────────────────────────────────────────
\section{Dirac Spinors as Time-Phase / Property Rotation}

A positive-energy Dirac spinor with momentum $p$ along $z$ (natural units):
\begin{equation}
  u(p,\uparrow) = N\begin{pmatrix}1\\0\\r\\0\end{pmatrix},\qquad
  u(p,\downarrow) = N\begin{pmatrix}0\\1\\0\\-r\end{pmatrix},\qquad
  r = \frac{p}{E+m},\quad N = \sqrt{\frac{E+m}{2E}}
\end{equation}

The \textbf{mixing angle} between temporal and spatial clock vectors:
\beq{
  \tan\frac{\theta_{\text{rel}}}{2} = \frac{|\mathbf{p}|}{E+m}
  \qquad\Longleftrightarrow\qquad
  \sin\theta_{\text{rel}} = \frac{v}{c}
}

\noindent Limits:
\[
  \theta_{\text{rel}} \xrightarrow{v/c\to 0} 0
  \quad\text{(NR: Pauli spinor sufficient)},
  \qquad
  \theta_{\text{rel}} \xrightarrow{v/c\to 1} \frac{\pi}{2}
  \quad\text{(massless: permanent orthogonality)}
\]

\subsection{Three-term decomposition of the Bell correlation}

\begin{equation}
  \boxed{E(a,b) =
    \underbrace{E_{LL}}_{\substack{\text{temporal}\\\times\,\text{temporal}}}
    +\underbrace{E_{SS}}_{\substack{\text{spatial}\\\times\,\text{spatial}}}
    +\underbrace{E_{LS}}_{\substack{\text{rotation}\\\text{coupling}}}
    = -\cos(a-b)}
\end{equation}

The NR clock model retains only $E_{LL}\approx -\tfrac{1}{2}\cos(a-b)$.
The Dirac small component ($E_{SS}+E_{LS}$) supplies the missing half,
verified numerically to machine precision ($<10^{-15}$) for all momenta.

% ─── 5 ────────────────────────────────────────────────────────────────────────
\section{Massless Particles -- Permanent Orthogonality}

For $m=0$: $\theta_{\text{rel}} = \pi/2$ always. The Dirac equation decouples:
\begin{align}
  i\sigma^\mu\partial_\mu\,\chi_R &= 0
  \qquad\text{(right-handed Weyl: spatial clock, }h=+\tfrac{1}{2}\text{)}\\
  i\bar\sigma^\mu\partial_\mu\,\chi_L &= 0
  \qquad\text{(left-handed Weyl: temporal clock, }h=-\tfrac{1}{2}\text{)}
\end{align}

The two clocks are permanently orthogonal and never mix.
This is why photons have exactly \textbf{2} polarization states rather than~4.

% ─── 6 ────────────────────────────────────────────────────────────────────────
\section{Higgs Field as Kuramoto Synchronizer}

The Yukawa interaction written in Weyl (clock) form:
\beq{
  \frac{d\varphi_L}{dt} = \omega + K\sin(\varphi_R - \varphi_L + \delta_{\text{CP}}),
  \qquad
  \frac{d\varphi_R}{dt} = \omega + K\sin(\varphi_L - \varphi_R - \delta_{\text{CP}})
}

where the \textbf{Kuramoto coupling equals the fermion mass}:
\begin{equation}
  \boxed{K = y_f\cdot\frac{v}{\sqrt{2}} = m_f}
\end{equation}
$y_f$ is the Yukawa coupling, $v=246\,\text{GeV}$ is the Higgs vacuum expectation value.
Without the Higgs ($K=0$) all fermions are massless and $\theta_{\text{rel}}=\pi/2$.

\subsection{Antiparticles as reversed clocks}

\begin{center}
\begin{tabular}{lcc}
\toprule
 & Particle $u(p)$ & Antiparticle $v(p)$ \\
\midrule
Phase & $e^{-iEt/\hbar}$ & $e^{+iEt/\hbar}$ \\
Clock direction & Forward ($+E$) & \textbf{Reversed} ($-E$) \\
Large/small components & Temporal dominant & Spatial dominant (swapped) \\
Sync equilibrium & $\varphi_L-\varphi_R\to+\delta_{\text{CP}}$ & $\varphi_L-\varphi_R\to-\delta_{\text{CP}}$\\
\bottomrule
\end{tabular}
\end{center}

Matter-antimatter asymmetry from CP-violating synchronization:
\begin{equation}
  \eta = \frac{N_{\text{matter}}-N_{\text{antimatter}}}{N_{\text{total}}}
  \approx \varepsilon\cdot\sin^2(\delta_{\text{CP}})
\end{equation}

% ─── 7 ────────────────────────────────────────────────────────────────────────
\section{Gravity as Bulk Clock Synchronization}

\subsection{Clock rate in a gravitational field}

\begin{equation}
  \omega(\mathbf{x})
  = \omega_\infty\sqrt{1+\frac{2\Phi(\mathbf{x})}{c^2}}
  \approx \omega_\infty\!\left(1+\frac{\Phi}{c^2}\right),
  \qquad \Phi = -\frac{GM}{r}
\end{equation}

\subsection{Poisson equation = Kuramoto field equation}

\begin{equation}
  \underbrace{\nabla^2\Phi = 4\pi G\rho}_{\text{Poisson (gravity)}}
  \qquad\longleftrightarrow\qquad
  \underbrace{\nabla^2\varphi_{\text{clock}} = -\rho\,K_{\text{grav}}}_{\text{Kuramoto (sync field)}}
\end{equation}

The gravitational potential \emph{is} the bulk clock synchronization field.
The clock-synchronization force recovers Newton exactly:
\begin{equation}
  a_{\text{clock}} = \frac{d\omega}{dr}\cdot\frac{c^2}{\omega_0} = \frac{GM}{r^2}
  \qquad\checkmark
\end{equation}

\subsection{Penrose objective reduction as gravitational clock decoherence}

Phase difference accumulated by superposition $|x_1\rangle+|x_2\rangle$:
\begin{equation}
  \delta\varphi(t) = \bigl[\omega(x_1)-\omega(x_2)\bigr]\cdot t
\end{equation}

Superposition coherence:
\begin{equation}
  \mathcal{C}(t) = \cos\!\left(\frac{\delta\varphi(t)}{2}\right)
\end{equation}

Collapse when $\mathcal{C}\to 0$ (i.e.\ $\delta\varphi=\pi$):
\beq{
  \tau_{\text{collapse}} = \frac{\pi\hbar}{E_G},
  \qquad
  E_G = \frac{Gm^2}{\Delta x}
  \label{eq:penrose}
}
Equation~\eqref{eq:penrose} is the Penrose formula, here derived from gravitational
clock decoherence. No observer is required.

% ─── 8 ────────────────────────────────────────────────────────────────────────
\section{Twistor Theory Connection}

Penrose's twistor $Z^\alpha=(\omega^A,\pi_{A'})$ maps to:
\[
  \omega^A \;\longleftrightarrow\; \psi_L\;\text{(temporal clock)},
  \qquad
  \pi_{A'} \;\longleftrightarrow\; \psi_R\;\text{(spatial clock)}
\]

Null twistor condition (massless):
\begin{equation}
  Z\cdot\bar Z = \omega^A\bar\pi_A + \bar\omega^{A'}\pi_{A'} = 0
  \qquad\Longleftrightarrow\qquad
  \theta_{\text{rel}} = 90^\circ
  \qquad\Longleftrightarrow\qquad
  m = 0
\end{equation}

Incidence relation (how spacetime position couples the two clocks):
\begin{align}
  \omega^A &= i\,x^{AA'}\pi_{A'}
  & &\text{(massless)}\\
  \omega^A &= i\,x^{AA'}\pi_{A'} + m\,\eta^A
  & &\text{(massive: Higgs coupling term)}
\end{align}
The massive incidence relation is the Dirac equation in Weyl form.
Gravity curves twistor space by deforming the incidence relation.

% ─── 9 ────────────────────────────────────────────────────────────────────────
\section{Temperature as Clock Phase Distribution Width}

\begin{equation}
  P(\varphi)
  \propto \exp\!\left(-\frac{\hbar\omega\,\varphi^2}{2k_BT}\right)
  \qquad\text{(classical limit }k_BT\gg\hbar\omega\text{)}
\end{equation}

Clock phase noise standard deviation:
\begin{equation}
  \sigma_\varphi(T) = \sqrt{\frac{k_BT}{\hbar\omega}},
  \qquad
  \sigma_\varphi(0) = \frac{1}{\sqrt{2}}\;\text{rad}
  \quad\text{(zero-point floor)}
\end{equation}

Bose-Einstein occupation number:
\begin{equation}
  \langle n(\omega)\rangle = \frac{1}{e^{\hbar\omega/k_BT}-1}
\end{equation}

Planck blackbody spectrum from thermal clock distribution:
\begin{equation}
  u(\omega,T) = \frac{\hbar\omega^3}{\pi^2 c^3}\cdot\frac{1}{e^{\hbar\omega/k_BT}-1}
  + \underbrace{\frac{\hbar\omega^3}{2\pi^2 c^3}}_{\text{zero-point residual}}
\end{equation}

Absolute zero is unreachable: the zero-point residual field
($\frac{1}{2}\hbar\omega$ per mode) is the minimum clock noise floor.

% ─── 10 ───────────────────────────────────────────────────────────────────────
\section{Brownian Motion from Clock Desynchronization}

At each molecular collision the momentum impulse is proportional to the clock phase
difference at contact:
\begin{equation}
  \delta p = K\sin(\varphi_i-\varphi_j) \approx K\,\delta\varphi,
  \qquad
  \delta\varphi \sim \mathcal{N}\!\left(0,\,\sigma_\varphi^2(T)\right)
\end{equation}

Mean squared impulse per collision:
\begin{equation}
  \langle(\delta p)^2\rangle = K^2\sigma_\varphi^2(T) = K^2\frac{k_BT}{\hbar\omega}
\end{equation}

The Langevin equation for a Brownian particle:
\begin{equation}
  m\ddot x = -\gamma\dot x + \xi(t),
  \qquad
  \langle\xi(t)\,\xi(t')\rangle = 2\gamma k_BT\,\delta(t-t')
\end{equation}

Diffusion coefficient predicted from clock parameters:
\beq{
  D_{\text{clock}}
  = \frac{K^2\,k_BT}{\hbar\omega}\cdot\frac{\tau_{\text{coll}}}{2m^2}
  \;\propto\; T
  \label{eq:diffusion}
}

Matching Eq.~\eqref{eq:diffusion} to the Stokes-Einstein relation $D=k_BT/(6\pi\eta r)$
\emph{fixes} the Kuramoto coupling:
\begin{equation}
  K = m\sqrt{\frac{\hbar\omega}{3\pi\eta r\,\tau_{\text{coll}}}}
\end{equation}

\subsection{New testable prediction}

Standard Stokes-Einstein has no dependence on molecular vibration frequency~$\omega$.
The clock model predicts:
\beq{
  D \;\propto\; \frac{1}{\sqrt{\omega}}
  \quad\text{for fixed }\eta,\,r,\,\tau
}
Measurable by single-particle tracking across solvents of different molecular mass.

% ─── 11 ───────────────────────────────────────────────────────────────────────
\section{Quantum-Classical Transition}

\subsection{Nelson's stochastic mechanics}

Nelson (1966) proved: Brownian motion in a background field with diffusion
$\nu=\hbar/2m$ satisfies the Schr\"odinger equation exactly.
The time-phase model identifies that background field as the zero-point residual wave field.

Total diffusion (quantum + thermal):
\begin{equation}
  \nu_{\text{total}}
  = \underbrace{\frac{\hbar}{2m}}_{\text{quantum (ZPF)}}
  + \underbrace{\frac{k_BT\,\tau}{m}}_{\text{thermal (Brownian)}}
\end{equation}

The quantum-classical crossover:
\begin{equation}
  \frac{k_BT\,\tau}{m} \gg \frac{\hbar}{2m}
  \;\;\Rightarrow\;\;\text{classical Brownian},
  \qquad
  \frac{k_BT\,\tau}{m} \ll \frac{\hbar}{2m}
  \;\;\Rightarrow\;\;\text{quantum (Schr\"odinger)}
\end{equation}

The Schr\"odinger's cat boundary coincides with the Penrose collapse time equalling
the thermal decoherence time -- two criteria derived from the same clock model.

\begin{equation}
  \tau_{\text{Penrose}} = \frac{\pi\hbar\,\Delta x}{Gm^2}
  \quad\Longleftrightarrow\quad
  \tau_{\text{thermal}} = \frac{m}{k_BT\,\tau_{\text{coll}}}
\end{equation}

\begin{center}
\begin{tabular}{lccc}
\toprule
System & $\hbar/2m$ (m$^2$/s) & $k_BT\tau/m$ at 300\,K & Regime \\
\midrule
Electron in atom ($\tau\sim 10^{-15}$\,s) & $3\times 10^{-5}$ & $10^{-21}$ & \textbf{Quantum} \\
Proton ($\tau\sim 10^{-13}$\,s) & $1.7\times 10^{-8}$ & $10^{-18}$ & Quantum \\
Protein ($10^{-22}$\,kg) & $5\times 10^{-13}$ & $10^{-14}$ & Borderline \\
1\,$\mu$m dust grain & $10^{-22}$ & $10^{-11}$ & \textbf{Classical} \\
\bottomrule
\end{tabular}
\end{center}

% ─── 12 ───────────────────────────────────────────────────────────────────────
\section{The Born Rule from Synchronization Statistics}

The Born rule --- $P = |\psi|^2$ --- is derived from Kuramoto synchronization
statistics rather than postulated.

\subsection{The Kuramoto order parameter}

For $N$ coupled oscillators with phases $\varphi_1,\ldots,\varphi_N$, the
synchronization amplitude is:
\begin{equation}
  R\cdot e^{i\Psi} = \frac{1}{N}\sum_{j=1}^{N} e^{i\varphi_j}
\end{equation}
The probability that a single oscillator locks to the collective is proportional
to~$R^2$ --- the squared modulus of the complex order parameter.

\subsection{Application to measurement}

A particle in superposition $\psi = \alpha|0\rangle + \beta|1\rangle$ arrives at
a detector with bulk phase~$\Phi_{\text{bulk}}$.  The synchronization probability
for each component depends on its complex coupling amplitude.  Averaging over the
uniformly distributed bulk phase eliminates cross terms:
\begin{equation}
  \langle e^{i(\varphi_0-\varphi_1)}\rangle_{\Phi_{\text{bulk}}} = 0
\end{equation}

Only the diagonal terms survive:
\beq{
  P(|0\rangle) = |\alpha|^2, \qquad P(|1\rangle) = |\beta|^2
  \label{eq:born}
}

\subsection{Why the squared modulus}

The wave function is complex because the internal clock is a phase oscillator
living on a circle in the complex plane.  The probability of synchronization is
the squared modulus because:
\begin{enumerate}[nosep]
  \item Coupling is linear in the complex amplitude (Kuramoto sine coupling)
  \item Probability is quadratic in the coupling (energy transfer $\propto$
        amplitude$^2$)
  \item The complex phase averages out over the bulk, leaving only
        $|\text{amplitude}|^2$
\end{enumerate}
The Born rule is the statement that quantum probability equals wave intensity ---
the natural measure for any physical wave coupling to a resonant detector.

% ─── 13 ───────────────────────────────────────────────────────────────────────
\section{Heisenberg Uncertainty from Clock Complementarity}

The uncertainty principle is not an axiom in this framework --- it is a geometric
consequence of the orthogonality between temporal and spatial clocks.

\subsection{Position-momentum uncertainty}

The temporal clock oscillates at frequency $\omega = E/\hbar$; the spatial clock
oscillates at wavenumber $k = p/\hbar$.  These are coupled by the incidence
relation (Section~8):
\begin{equation}
  \omega^A = i\,x^{AA'}\pi_{A'} + m\,\eta^A
\end{equation}

A sharp position measurement locks the spatial clock phase~$\varphi_R$ to a
definite value at a point~$x$.  But the incidence relation couples $\varphi_R$
to~$\varphi_L$ through the spacetime position --- locking one forces the other
into a superposition of frequencies.  This is the Fourier relationship between
position and momentum representations, here derived from clock geometry rather
than postulated.

The mixing angle~$\theta_{\text{rel}}$ encodes the tradeoff.
When $\theta_{\text{rel}}\to 0$ (rest frame), the temporal clock dominates and
energy is sharp --- but position is maximally uncertain.
When $\theta_{\text{rel}}\to\pi/2$ (ultrarelativistic), the spatial clock
dominates and momentum is sharp --- but energy-time is maximally uncertain.

The minimum uncertainty product follows from the Kuramoto coupling between the
two clocks.  The synchronization force
\begin{equation}
  \frac{d\varphi_L}{dt} = \omega + K\sin(\varphi_R - \varphi_L)
\end{equation}
ensures that locking~$\varphi_L$ (energy measurement) introduces noise
in~$\varphi_R$ (momentum) of order
$\delta\varphi_R \geq 1/(2\,\delta\varphi_L)$, because the sine coupling
transfers phase uncertainty bidirectionally.  Converting to physical units
($\Delta E = \hbar\,\Delta\omega$, $\Delta p = \hbar\,\Delta k$):
\beq{
  \Delta x \cdot \Delta p \;\geq\; \frac{\hbar}{2}
  \label{eq:hup-xp}
}

\subsection{Energy-time uncertainty}

The particle clock runs at $\omega = E/\hbar$.  A measurement lasting
time~$\Delta t$ samples $\Delta t/(2\pi/\omega)$ oscillation cycles.  The
frequency resolution of any oscillator observed for time~$\Delta t$ is
$\Delta\omega \geq 1/(2\,\Delta t)$.  Multiplying by~$\hbar$:
\beq{
  \Delta E \cdot \Delta t \;\geq\; \frac{\hbar}{2}
  \label{eq:hup-et}
}
This is identical to the Fourier bandwidth theorem --- but here the ``signal''
is a physical clock, not an abstract wave.

\subsection{Zero-point phase noise as the uncertainty floor}

From Section~9, the zero-point phase noise at $T=0$ is
$\sigma_\varphi(0) = 1/\sqrt{2}$~rad.  This residual prevents perfect phase
locking of either clock.  The zero-point energy $\frac{1}{2}\hbar\omega$ per
mode is the energy cost of this minimum phase uncertainty.  The uncertainty
principle is thus the statement that the zero-point field cannot be removed ---
complete synchronization of both clocks simultaneously is forbidden by the
geometry of spinor space.

\subsection{Nyquist--Shannon bandwidth limit}

The connection between the zero-point floor and the uncertainty principle can
be made sharper through a Nyquist--Shannon argument.  A Dirac particle's
internal clock oscillates at the Compton frequency
$\omega_C = mc^2/\hbar$.  By the Nyquist sampling theorem, this clock has a
maximum spatial bandwidth: it can resolve phase structures no finer than the
Compton wavelength $\lambda_C = h/(mc)$.  This sets an irreducible minimum
position uncertainty:
\beq{
  \Delta x_{\min} \sim \frac{\hbar}{mc}
  \label{eq:nyquist-dx}
}
For a particle with momentum spread~$\Delta p$, the de~Broglie relation gives
a spatial frequency $k = p/\hbar$ that the clock must resolve.  The Nyquist
limit requires the clock bandwidth $mc/\hbar$ to exceed~$k$, yielding:
\beq{
  \Delta x \cdot \Delta p \;\geq\; \frac{\hbar}{2}
  \label{eq:nyquist-hup}
}
The uncertainty principle is thus a \textbf{bandwidth limitation of the
particle's own internal clock}.  A heavier particle has a faster clock and
therefore finer spatial resolution, but the resolution is always finite.  No
oscillator can resolve frequencies beyond its own Nyquist limit.

\subsection{The commutator from clock coupling}

The canonical commutation relation
\begin{equation}
  [\hat x,\,\hat p\,] = i\hbar
\end{equation}
emerges from the non-commutativity of sequential clock measurements.  Measuring
position (locking~$\varphi_R$, then reading~$\varphi_L$) and measuring momentum
(locking~$\varphi_L$, then reading~$\varphi_R$) do not commute because the
Kuramoto sine coupling is nonlinear --- the order of synchronization matters.
The imaginary unit~$i$ reflects the $\pi/2$ phase shift between the two clocks
at the point of maximum complementarity ($\theta_{\text{rel}}=\pi/2$).

% ─── 14 ───────────────────────────────────────────────────────────────────────
\section{Complete Dictionary}

\begin{center}
\begin{tabular}{p{4.2cm}p{3.8cm}p{4.2cm}}
\toprule
\textbf{Time-phase model} & \textbf{Standard physics} & \textbf{Geometric (Twistor)} \\
\midrule
Temporal clock $\psi_L$ & Left Weyl spinor & $\omega^A$ \\
Spatial clock $\psi_R$ & Right Weyl spinor & $\pi_{A'}$ \\
L-R coupling $= m$ & Dirac mass term & Non-null: $Z\cdot\bar Z\neq 0$ \\
$\theta_{\text{rel}}=90^\circ$ & Massless particle & Null twistor \\
Higgs synchronizer & Yukawa coupling & Twistor cohomology \\
Gravity $=\nabla^2\varphi$ & Newtonian potential & Twistor space curvature \\
Penrose collapse & Wavefunction collapse & $\delta\varphi=\pi$ \\
Temperature & $k_BT$ & $\sigma_\varphi^2 = k_BT/\hbar\omega$ \\
Brownian motion & Stokes-Einstein & $D=K^2\sigma_\varphi^2\tau/2m^2$ \\
Zero-point field & $\tfrac{1}{2}\hbar\omega$ per mode & $\sigma_\varphi(0)=1/\sqrt{2}$ \\
Antiparticle & $e^{+iEt/\hbar}$ & Reversed clock \\
Matter asymmetry & Baryogenesis & $\eta\approx\varepsilon\sin^2\delta_{\text{CP}}$ \\
Born rule $P=|\psi|^2$ & Postulated & Sync probability $=|R|^2$ \\
Uncertainty principle & $\Delta x\,\Delta p\geq\hbar/2$ & Clock orthogonality: $\delta\varphi_L\,\delta\varphi_R\geq\tfrac{1}{2}$ \\
Measurement & Wavefunction collapse & Kuramoto re-sync to bulk \\
Entanglement & Non-separable state & Synchronized clocks: $\varphi_A=\varphi_B$ \\
Bell's theorem & No local HV model & $\Phi_{\text{bulk}}$ shared $\Rightarrow$ factorization fails \\
Quantum-classical & Decoherence & $k_BT\tau/m \gg \hbar/2m$ \\
Nelson stochastic QM & $\nu=\hbar/2m$ & Zero-point field diffusion \\
\bottomrule
\end{tabular}
\end{center}

\vspace{2em}
\hrule
\vspace{1em}
\small\textit{%
Generated from conversations between a radiologist and Claude~(Anthropic).
Mathematical framework developed by AI\@; physical intuitions by the author.
Inspirations: de~Broglie (pilot wave), Bohm (holographic order),
Penrose (twistors, OR, CCC), Kuramoto (phase synchronization), Nelson (stochastic QM).%
}

\end{document}
